% !TeX root = ../dp-puzzle_zh.tex
\section{背景与动机}
\label{sec:background}

\subsection{浮点精度与微缩放格式}
\label{sec:formats}

以 FP16 为例，二进制浮点编码包含符号、指数和小数字段（fraction field）。本文将小数字段与隐含前导位组成的无符号定点数称为尾数（mantissa）。正规数的前导位为 1，非正规数和零的前导位为 0。指数位宽决定动态范围，小数字段位宽决定相对分辨率。TF32 保留 FP32 的指数范围并缩短尾数，以降低张量计算开销；FP16 与 BF16 则分别侧重分辨率与动态范围~\cite{burgess2019bfloat16}。开放计算项目（Open Compute Project, OCP）定义的微缩放格式（Microscaling, MX）~\cite{ocpmx2023}将元素表示为 $x_{j,i}=X_jP_{j,i}$，其中 $X_j$ 为组级共享缩放因子（scaling factor），$P_{j,i}$ 为元素解码值。MXFP8、MXFP6 与 MXFP4 分别采用 E5M2/E4M3、E3M2/E2M3 与 E2M1 元素编码，均以 32 个元素为一组共享一个 E8M0 缩放因子，以扩展可表示的数值范围。

然而，权重与激活同时量化为 MXFP4 时仍有明显的精度损失。其 E8M0 共享缩放因子只能取 2 的幂，是量化误差的主要来源~\cite{egiazarian2026bridging,zhang2026benchmarking}。为改善量化精度，NVFP4 保留 E2M1 元素编码，将分组缩小至 16，并采用 E4M3 组级与 FP32 张量级两级缩放，提升表示精度~\cite{nvidia2024nvfp4}；HiF4 则采用 S1P2 定点元素与分层微指数缩放，兼顾组内动态范围与分辨率~\cite{hifloat4}。NVFP4 和 HiF4 的块缩放因子含小数字段，施加块缩放时还需在指数调整之外进行缩放因子的尾数乘法。

\subsection{面向微缩放格式的点积计算}
\label{sec:tensor_organization}

常规数据格式与微缩放格式~\cite{ocpmx2023}的点积累加可统一表示为
\begin{equation}
\begin{aligned}
\operatorname{DP}_K(\mathbf{a},\mathbf{b},c)
&=c+\sum_{k=0}^{K-1}a_kb_k \\
&=c+\sum_{j=0}^{G_s-1}
\widetilde{X}_j^{(A)}\widetilde{X}_j^{(B)}
\underbrace{\sum_{i=0}^{g-1}
P_{j,i}^{(A)}P_{j,i}^{(B)}}_{q_j},
\end{aligned}
\label{eq:blockdot}
\end{equation}
其中，两路输入沿归约维度采用相同的缩放分组边界，$g$ 为分组大小，$G_s=\lceil K/g\rceil$，末块补零；$q_j$ 为块内点积，$\widetilde{X}_j$ 为等效块缩放因子。常规浮点格式取 $\widetilde{X}_j=1$，$P_{j,i}$ 为原始元素值；MX 格式取 $\widetilde{X}_j=X_j$；HiF4 的局部微指数计入 $P_{j,i}$。

\subsection{多精度点积单元设计与设计空间探索相关工作}
\label{sec:datapath_sharing}

在乘法复用方面，Bit Fusion 以移位加法组合低位宽乘法单元；Park 等人在 FP16 尾数乘法与 INT8/INT4 点积之间共享乘法器与归约通路；Abdel-Aziz 等人通过时间分解支持更高精度，并分析了浮点指数对齐与加法的额外开销~\cite{bitfusion2018,parkMultimode8KMACHWutilizationaware2022,abdelaziz2021overheads}。在多格式计算与融合归约方面，FPnew 以可配置的格式切片支持多种浮点运算；MiniFloats 将两项低精度乘积与高精度累加项融合，减少重复的规格化与舍入~\cite{mach2020fpnew,bertaccini2022minifloat}。TransDot 复用乘法部分积与可分区移位器；Ten-Four 将多种格式的点积接入共享归约后端；MXDOTP 采用统一的内部元素表示，在点积通路中融合 MX 缩放并累加至 FP32~\cite{transdot2026,rout2026tenfourv3,islamoglu2025mxdotp}。
在 $C$ 累加阶段的选择上，乘积既可与累加项共同对齐归约，也可在形成缩放点积后送入浮点加法器；接入位置决定对齐范围与累加通路开销~\cite{bertaccini2022minifloat,lutzFusedFP84Way2024}。此外，内部累加顺序、截断位置与舍入规则共同决定数值结果~\cite{sun2023tensorcores,xie2025mmasim}。在设计空间探索方面，Crafton 等与 MX-Like 进一步分析了低精度格式、计算组织与硬件成本之间的权衡~\cite{paretomac2025,zhao2026mxlike}，但均未涉及算力配比的选择问题。

表~\ref{tab:related_features} 对比了本文与相关点积单元工作的精度格式覆盖范围。与上述工作相比，本文在支持多种精度格式的基础上，通过增加低精度专用单元改变点积单元内各格式间的峰值算力，从而实现多种算力配比。

\subsection{算力配比的设计动机}
\label{sec:commercial_ratios}
\label{sec:ratio_motivation}

表~\ref{tab:vendor_ratios} 列出了当前主流 GPGPU 与 NPU 厂商的最新产品与研发路线，体现了工业界对 $1:2:6$、$1:4:8$ 与 $1:2:8$ 三种配比的认同，但均未说明其选择依据，且既有选择仍局限于较小的探索空间。因此，本文进一步考察算力配比的可选组合，将 $1:4:4$、$1:2:12$、$1:2:16$、$1:6:12$ 与 $1:8:16$ 纳入探索，与上述四种配比（含基准 $1:2:4$）共同构成九种候选设计空间。本文希望在未来负载精度格式迁移的背景下，回答如何在硬件设计中权衡开销与收益这一问题。